% Propietario conservado: /Users/ruben/Documents/New project/output/REVISION_ARGUMENTAL_APERTURAS_CIERRES_20260915/VII/delivery/MANUSCRIPT_VII_ES_EN_COMPLETE/source_es/manuscrito/50_dilucion_y_rebote.tex
% Artículo VII: pruebas preexistentes incorporadas al cuerpo.
% Propietario íntegro: fuentes_conservadas/17_rebote_einstein_cartan.tex.
% Incluye la ley de escala y las pruebas de rebote y persistencia.
% La amplitud s_0 es la salida de la realización material antecedente;
% este fragmento no certifica por sí solo la evaluación de esa realización.

\section{Dilution de la mémoire et persistance du rebond}
\label{vii:vii:sec:dilucion-rebote}

L'évolution de la mémoire et l'expansion spatiale sont deux lectures d'une même suite d'états. L'indice de raffinement peut être éliminé entre les deux ; cette élimination détermine l'exposant de dilution que la réalisation cosmologique transmet à sa contribution torsionnelle. L'amplitude initiale dépend de l'évaluation du courant matériel concret dans sa réponse variationnelle. La loi d'échelle et le théorème de rebond qui suivent conservent cette dépendance.

Le théorème~\ref{vii:vii:thm:amplitud-material-explicita} détermine cette
amplitude comme fonctionnelle de l'état matériel et du volume comobile.
Le corollaire~\ref{vii:vii:cor:dominio-positivo-conservado} fournit
un domaine explicite d'états d'amplitude positive et conservée.
Dans ce domaine, l'évaluation de \(s_0\) précède l'intégration
cosmologique développée ci-dessous. La densité ordinaire
\(\rho_0>0\) et sa loi barotropique sont spécifiées en outre dans la
section~\ref{vii:vii:sec:umbral-rebote} : le témoin de positivité de l'
amplitude n'évalue pas par lui-même cette densité matérielle.

\subsection{Loi de dilution \(a^{-6}\)}
\label{vii:vii:sec:dilucion}

\begin{proposition}[Élimination de l'indice de mémoire]
\label{vii:vii:prop:dilucion}
Soient \(a_{\mathrm{ref}}>0\), \(\varphi>1\) et les lecteurs de la route
\[
 \frac{a_n}{a_{\mathrm{ref}}}=\varphi^{n/3},\qquad
 M_n=\varphi^{-2n}.
\]
Alors
\[
 M_n=\left(\frac{a_n}{a_{\mathrm{ref}}}\right)^{-6}
\]
pour tout indice de la suite. Si la réalisation de la densité
quadratique est \(s_n^2=s_*^2M_n\), avec \(s_*^2>0\), on obtient
\[
 s_n^2=s_*^2\left(\frac{a_n}{a_{\mathrm{ref}}}\right)^{-6}.
\]
\end{proposition}

\begin{proof}
En élevant la première égalité à la puissance \(-6\), on obtient \(\varphi^{-2n}\), qui est exactement le second lecteur. Multiplier par l'amplitude initiale reproduit la seconde affirmation. L'exposant est déterminé par le quotient des exposants des deux lecteurs ; l'amplitude est l'évaluation de l'état de référence.
\end{proof}

Dans la réalisation homogène continue, cette loi constitutive est conservée :
la densité torsionnelle s'écrit \(s_0a^{-6}\), où la normalisation
du facteur \(a\), le coefficient matériel et l'amplitude de référence
ont été absorbés dans \(s_0\). La loi d'échelle et l'évaluation de
\(s_0\) sont des opérations successives de la construction.

\subsection{Seuil homogène et transversalité}
\label{vii:vii:sec:umbral-rebote}

Considérons la carte homogène, spatialement plate, en unités \(c=1\),
avec \(G>0\) constant, des densités positives \(\rho_0,s_0\) et un paramètre
barotrope constant \(w<1\). La réalisation spécifiée satisfait
\(\Lambda_{\mathrm{eff}}=0\) et les équations
\begin{align}
 \rho_{\mathrm m}(a)&=\rho_0a^{-3(1+w)},&
 \rho_{\mathrm s}(a)&=s_0a^{-6},\\
 H^2&=\frac{8\pi G}{3}
        \bigl(\rho_{\mathrm m}-\rho_{\mathrm s}\bigr)
        =:F_0(a),&
 \dot H&=-4\pi G
        \bigl((1+w)\rho_{\mathrm m}-2\rho_{\mathrm s}\bigr).
 \label{vii:vii:eq:dinamica-homogenea}
\end{align}
Les couplages appartiennent à la section dimensionnelle déjà construite.
L'expression de \(\rho_{\mathrm s}\) intègre le signe de sa contribution
dans cette réalisation gravitationnelle.

\begin{theorem}[Existence et unicité du seuil transversal]
\label{vii:vii:thm:rebote}
Il existe un unique \(a_{\mathrm b}>0\) tel que \(F_0(a_{\mathrm b})=0\) :
\[
 a_{\mathrm b}^{\,3(1-w)}=\frac{s_0}{\rho_0}.
\]
La région permise par \(H^2\geq0\) est \(a\geq a_{\mathrm b}\).
Au seuil, avec
\(\rho_{\mathrm b}=\rho_{\mathrm m}(a_{\mathrm b})
                  =\rho_{\mathrm s}(a_{\mathrm b})>0\),
\[
 \dot H_{\mathrm b}=4\pi G(1-w)\rho_{\mathrm b}>0.
\]
La solution qui traverse \(H=0\) présente un minimum strict de \(a\), avec le développement local
\[
 a(t)=a_{\mathrm b}\left[
 1+2\pi G(1-w)\rho_{\mathrm b}(t-t_{\mathrm b})^2
 +O((t-t_{\mathrm b})^3)\right].
\]
\end{theorem}

\begin{proof}
La factorisation
\[
 F_0(a)=\frac{8\pi G}{3}a^{-6}
             \bigl(\rho_0a^{3(1-w)}-s_0\bigr)
\]
réduit ses zéros à celui d'une fonction strictement croissante de
\((0,\infty)\) sur \((-s_0,\infty)\). Cela démontre l'existence,
l'unicité et le signe de \(F_0\) de part et d'autre du seuil.
La seconde équation de \eqref{vii:vii:eq:dinamica-homogenea}, évaluée
à l'égalité des densités, donne la formule de \(\dot H_{\mathrm b}\).
Comme \(\dot a=aH\), au seuil
\(\ddot a_{\mathrm b}=a_{\mathrm b}\dot H_{\mathrm b}>0\).
Les équations sont lisses pour \(a>0\), de sorte que leur développement de
Taylor donne le développement indiqué. Le changement de signe de \(H\) est
transverse : il passe de la contraction à l'expansion.
\end{proof}

Pour \(w=1\), les deux termes ont le même exposant et l'égalité des
amplitudes détermine un cas dégénéré. Pour \(w>1\), le signe de
\(\dot H\) lors d'une coïncidence des densités est négatif.
Le domaine \(w<1\) du théorème exprime exactement l'ordre entre les
deux exposants de dilution.

\subsection{Persistance locale sous perturbations contrôlées}
\label{vii:vii:sec:persistencia}

La persistance du seuil utilise deux estimations indépendantes : le contrôle du signe aux extrémités et celui de la dérivée dans l'intervalle. Soit \(I=[a_-,a_+]\subset(0,\infty)\), avec \(a_-<a_{\mathrm b}<a_+\), choisi de sorte que
\[
 F_0(a_-)<0<F_0(a_+),\qquad
 m_I:=\inf_{a\in I}F_0'(a)>0.
\]
Un tel intervalle existe parce que
\[
 F_0'(a_{\mathrm b})
   =\frac{8\pi G}{a_{\mathrm b}}(1-w)\rho_{\mathrm b}>0.
\]
Nous définissons
\(\delta_I=\min\{-F_0(a_-),F_0(a_+)\}>0\).

\begin{theorem}[Persistance locale du rebond transversal]
\label{vii:vii:thm:persistencia}
Soit \(R\in C^1(I)\) une perturbation qui satisfait
\[
 \|R\|_{\infty,I}<\delta_I,\qquad
 \|R'\|_{\infty,I}<m_I.
\]
Alors \(F_R=F_0+R\) possède un unique zéro \(a_R\in(a_-,a_+)\) et \(F_R'(a_R)>0\). La réalisation \(H^2=F_R(a)\) admet le rebond local transversal, avec
\[
 \dot H_R=\frac{a_R}{2}F_R'(a_R)>0
\]
au seuil. Si \(R(a,\eta)\) est conjointement \(C^1\), les zéros
persistants dépendent localement de manière \(C^1\) du paramètre \(\eta\).
\end{theorem}

\begin{proof}
La première borne conserve \(F_R(a_-)<0<F_R(a_+)\) ; le théorème des valeurs intermédiaires fournit un zéro intérieur. La seconde donne \(F_R'(a)\geq m_I-\|R'\|_{\infty,I}>0\) dans \(I\), de sorte que le zéro est unique et simple.

En \(a>a_R\) proche du zéro,
\(F_R(a)=F_R'(a_R)(a-a_R)+o(a-a_R)\).
L'intégrale
\[
 t-t_{\mathrm b}=\int_{a_R}^{a}
          \frac{d\xi}{\xi\sqrt{F_R(\xi)}}
\]
est finie et strictement croissante. Son inverse définit la branche
en expansion ; la réflexion temporelle définit la branche en contraction.
Toutes deux se rejoignent avec \(\dot a(t_{\mathrm b})=0\) et
\[
 a(t)-a_R
 =\frac{a_R^2F_R'(a_R)}4(t-t_{\mathrm b})^2
   +o((t-t_{\mathrm b})^2).
\]
Hors du seuil, la dérivation de \(H^2=F_R(a)\) donne
\(\dot H=aF_R'(a)/2\) ; sa limite continue donne l'expression positive
en \(a_R\). La régularité \(C^1\) de \(F_R\) garantit la continuité
de l'accélération à cette jonction. Enfin,
\(\partial_aF_R(a_R)>0\) permet d'appliquer le théorème des fonctions
implicites lorsque la perturbation dépend conjointement de \((a,\eta)\).
\end{proof}

La première borne contrôle l'existence ; la seconde assure l'unicité locale et la transversalité. La stabilité démontrée correspond à ce seuil dans l'intervalle sélectionné. L'amplitude torsionnelle intervient dans sa localisation, tandis que la relation entre exposants détermine l'orientation temporelle du franchissement.
